\sgaStableTarget{sga6.e07.remark-3-8}\subsection*{Remark 3.8.}

When @@M0@@ is noetherian and @@M1@@ regular, 3.7 determines the abelian group @@M2@@ entirely from the @@M3@@-rings @@M4@@, @@M5@@, the maps @@M6@@ between these rings, and the element @@M7@@ of @@M8@@: indeed, @@M9@@ can be computed from @@M10@@ and @@M11@@ by (VI 1.1). We shall see in the next section that the multiplicative structure of @@M12@@ is determined by the same data, as is the @@M13@@-structure after tensoring with @@M14@@ (this restriction is probably unnecessary: see 4.5).

\sgaStableTarget{sga6.e07.sec.4}\section*{4. Regular immersions and the filtration of @@M15@@}

Throughout this section we keep the notation and hypotheses of 3.1.

\sgaStableTarget{sga6.e07.lem.4-1}\textbf{Lemma 4.1.} \uline{Suppose that there exists on} @@M16@@ \uline{a locally free sheaf of finite type} @@M17@@, \uline{such that}, \uline{if} @@M18@@ \uline{are the classes in} @@M19@@ \uline{of} @@M20@@, \uline{we have} @@M21@@. \uline{Then the following congruence holds}:
@@M22@@

Let @@M23@@ be the rank of @@M24@@ (which we may suppose constant); the rank of @@M25@@ is then @@M26@@. By definition:
@@M27@@
Now @@M28@@. Thus for every @@M29@@, @@M30@@ is divisible by @@M31@@. Consequently:
@@M32@@
